
% Definitionen/Pakete für Handout und Aufgaben

\KOMAoptions{parskip=half, headinclude=true, headsepline=true}

% Pakete
\usepackage{fontspec}
\setmainfont{Latin Modern Roman}
\usepackage[ngerman]{babel}
\usepackage{csquotes}
\usepackage{graphicx}
\newcommand*{\imagesourcefolder}{./Bilder}
\graphicspath{{\imagesourcefolder/}}
\usepackage{xcolor}
\usepackage{wrapfig}
\usepackage{mathtools}
\usepackage[per-mode=fraction]{siunitx}
\usepackage{tikz,pgfplots}
\usepgfplotslibrary{polar}
\usetikzlibrary{arrows.meta}
\usetikzlibrary{plotmarks}
\usetikzlibrary{calc}
\pgfplotsset{compat=newest}
\usepackage{tabularx}
\usepackage{booktabs}

% Layout
\usepackage[
origlayout=true, automark, colors={mat},
headsepline=4pt,
intern  % spare Tinte
]{URpagestyles}
\pagestyle{URheadings}
\renewcommand{\titlepagestyle}{URtitle}
\addtokomafont{section}{\Large\normalfont\bfseries}

% Lückentext/Antworten
\usepackage[thickness=0pt,textcolor=gray]{cloze}
\clozesetfont{\LARGE\itshape}
\makeatletter
\newif\ifhide
\@ifclasswith{scrartcl}{hide}{\hidetrue}{\hidefalse}
\makeatother
\newcommand{\Loesung}[1]{\ifhide\relax\else{\clozefont{}\color{gray}#1}\fi}
\newcommand{\ausfuelltext}[1]{\cloze{#1\strut}}
\newcommand{\formel}[1]{\fcolorbox{gray}{white}{#1}}
\newcommand{\canvas}[3][c]{
  \parbox[#1]{#2}{~\\[-\baselineskip]%
    \formel{#3}%
    \\\hbox{}
  }
}

% Titel, Metadaten
\usepackage[colorlinks=true, linkcolor=blue,
pdfauthor={Gesina Schwalbe}
]{hyperref}
%\subject{Schülerzirkel am XX.XX.2018}
\date{}
\author{Gesina Schwalbe}

\clozesetfont{\itshape}

% Titel
\hypersetup{pdftitle={Aufgaben: Vom Skalarprodukt zur
    Bildbearbeitung}}
\title{Aufgaben}
\subtitle{zum Vortrag \enquote{Vom Skalarprodukt zur Bildbearbeitung}}

\begin{document}
\maketitle

\section*{Eigenschaften von Vektoren}
\begin{enumerate}
\item \emph{(Distributivität von Vektoraddition und Skalarmultiplikation)}
  Berechne folgende Vektoren:
  \begin{enumerate}
  \item $2 \cdot \left[ (2,1) + (-3,1) \right]
    \text{\cloze{$=(-2,4)$}}$
    \\$\left[ 2\cdot(2,1)\right] + \left[ 2\cdot(-3,1) \right]
    \text{\cloze{$=(-2,4)$}}$
  \item $-3 \cdot \left[ (1,3,-2) + (-1, -2, 3) \right]
    \text{\cloze{$=(0, -3, -3)$}}$
    \\$\left[ -3\cdot(1,3,-2) \right] + \left[ -3\cdot(-1, -2, 3)\right]
    \text{\cloze{$=(0, -3, -3)$}}$
  \end{enumerate}
  Kannst du folgende allgemeine Regel für Vektoren
  $x=(x_1,\dotsc,x_n)$ und $y=(y_1,\dotsc,y_n)$ bestätigen?
  \begin{gather*}
    \SwapAboveDisplaySkip
    a \cdot \left[ x + y \right]
    = \left[ a\cdot x \right]
    + \left[ a\cdot y \right]
  \end{gather*}
\item \emph{(Bilinearität des Skalarprodukts I)}
  Berechne folgende Skalarprodukte:
  \begin{enumerate}
  \item $\left[ (1, -2) + (-3, 1) \right] \circ (1, -3)
    \text{\cloze{$=(-2,-1)\circ(1,-3)= 1$}}$
    \\$\left[(1,-2)\circ(1,-3)\right] + \left[(-3,1)\circ(1,-3)\right]
    \text{\cloze{$=(-2,-1)\circ(1,-3)= 1$}}$
  \item $(3, 2) \circ \left[ (-5, 3) + (4, -3)\right]
    \text{\cloze{$=(3,2)\circ(-1,0)= -3$}}$
    \\$\left[(3,2)\circ(-5,3)\right] + \left[(3,2)\circ(4,-3)\right]
    \text{\cloze{$=(3,2)\circ(-1,0)= -3$}}$
  \end{enumerate}
  Kannst du folgende allgemeinen Regeln für Vektoren
  $x=(x_1,\dotsc,x_n)$, $y=(y_1,\dotsc,y_n)$, und $z=(z_1,\dotsc,z_n)$
  bestätigen?
  \begin{gather*}
    \left[x + y\right] \circ z
    = \left[x\circ z\right] + \left[y\circ z\right]
    \\
    x \circ \left[y + z\right]
    = \left[x\circ y\right] + \left[x\circ z\right]
  \end{gather*}
\item \emph{(Bonus: Bilinearität des Skalarprodukts II)}
  Berechne folgende Skalarprodukte:
  \begin{enumerate}
  \item $3 \cdot \left[ (-1, 2) \circ (3, 1) \right]
    \text{\cloze{$=3 \cdot -1= -3$}}$
  \item $\left[ 3\cdot(-1, 2) \right] \circ (3, 1)
    \text{\cloze{$=3 \cdot -1= -3$}}$
  \item $(-1, 2) \circ \left[ 3\cdot (3, 1)\right]
    \text{\cloze{$=3 \cdot -1= -3$}}$
  \end{enumerate}
  Kannst du folgende allgemeine Regel für Vektoren
  $x=(x_1,\dotsc,x_n)$ und $y=(y_1,\dotsc,y_n)$ bestätigen?
  \begin{gather*}
    a \cdot \left[ x \circ y \right]
    = \left[ a\cdot x \right] \circ y
    = x \circ \left[ a\cdot y \right]
  \end{gather*}
\item \emph{(Bonus: Geschwindigkeitsvektoren)}
  Wähle ein Koordinatensystem und zeichne den Geschwindigkeitsvektor
  eines Steins ein, der mit \SI{10}{\meter\per\second} (Länge unseres
  Geschwindigkeitsvektors) im nicht ganz optimalen Winkel von
  \SI{30}{\degree} auf eine Wasseroberfläche trifft.
  Welchen senkrechten Geschwindigkeitsanteil hat unser Stein?
  \Loesung{Der Vektor ist ca. $(9,66 , 2,59)$.}
  \par
  \emph{(Bonus-Bonus)} Weißt du, warum er, wenn richtig geworfen, nicht einsinkt?
% \item \emph{Bonus: Kraftvektoren}
%   Eine Lampe hängt an zwei gleich langen Seilen auf, die beide im
%   \SI{45}{\degree}-Winkel von gegenüberliegenden Wänden hängen. 
%   Die Lampe ist ordentlich schwer und zieht mit einer Kraft von
%   10~Newton (das ist die Länge unseres Kraftvektors), also
%   10~Schokoladentafeln, senkrecht nach unten.
%   Damit sie nicht herunterfällt, müssen jeweils in Richtung der beiden
%   Seile zwei so starke Kräfte ziehen, dass ihr gemeinsamer,
%   d.\,h. addierter, senkrechter Anteil die 10~Newton aufhebt
%   (also ein -10~Newton langer Vektor in senkrechter Richtung).
%   \begin{enumerate}
%   \item Wähle ein Koordinatensystem für Kräfte und zeichne die 10
%     bzw. -10~Newton sowie die Richtung der beiden Fadenkräfte ein.
%   \item Finde die Fadenkräftevektoren und zeichne sie ein. Wie lang
%     (d.\,h. stark) sind diese?
%   \item Was ist der horizontale Anteil eines Fadenvektors?
%     (Das ist dann die Kraft, die die Wand aushalten muss)
%   \item[] \emph{(Bonus)} Probiere das ganze nochmal mit gespannteren
%     Seilen, z.\,B. mit einem Winkel von \SI{170}{\degree} zwischen den
%     Seilen. Wie stark sind dann die horizontalen Kräfte, die die Wand
%     aushalten muss?
%     Traust du dich dann noch, unter der Lampe durchzulaufen?
%   \end{enumerate}
\end{enumerate}

\section*{Bildbearbeitung mit Gimp in Theorie und in der Praxis}
\begin{enumerate}
\item Welcher Vergleichsbildausschnitt verändert bei einer Faltung ein
  Bild gar nicht?
  \Loesung{Der mittlere Eintrag 1, sonst nur 0}
  \par
  Mit welchem Vergleichsbildausschnitt kannst du das Bild mit einer
  Faltung einfach nur heller, und mit welchem einfach nur dunkler
  machen? Was hat er mit dem Ergebnis von oben zu tun? 
  \Loesung{Der Faktor ist der mittlere Eintrag, alle anderen sind
    0. D.h. es ist ein Vielfaches des Kerns von oben.}
\item Welcher Vergleichsbildausschnitt könnte diese Faltung hervorrufen?
  \begin{center}
    \parbox[c]{0.4\linewidth}{
      \includegraphics[width=\linewidth]{horizontale_kante-dunkel-original}
      \\\hbox{}}
    $\Rightarrow$
    \parbox[c]{0.4\linewidth}{
      \includegraphics[width=\linewidth]{horizontale_kante-dunkel}
      \\\hbox{}}
  \end{center}
  \Loesung{Mögliche Antworten:
    \begin{gather*}
      \begin{bmatrix}
        0 &-1 &0\\
        0 & 1 &0\\
        0 & 0 &0
      \end{bmatrix}
      \qquad
      \begin{bmatrix}
        -1 &-1 &-1\\
         1 & 1 & 1\\
         0 & 0 & 0
      \end{bmatrix}
    \end{gather*}}
\item Überlege dir, wie ein Kernel für folgende Probleme aussehen könnte,
  probiere deine Ideen zusammen mit deiner Gruppe aus und verbessere
  sie:
  \begin{description}
  \item[Künstliche Unschärfe] Ersetze ein Pixel durch den
    Durchschnittswert seiner umliegenden Pixel.
    Inwiefern stellt der Vergleichsbildausschnitt eine Unschärfe dar?
    \Loesung{Am einfachsten: Setze alle 9 Kernwerte auf 1 und den
      Divisor auf $\frac 19$}
    \\Bonus: Was heißt unscharf bei Bildern für benachbarte
    Pixelwerte?
    \Loesung{Die Helligkeit benachbarter Pixel unterscheidet sich
      wenig – alles Einheitsbrei.}
  \item[Ecken-/Kantenerkennung] Markiere alle Kantenpixel hell, die
    übrigen dunkel.
    Tipp: Finde erst Vergleichsbildausschnitte für verschiedene
    Kantentypen und addiere sie.
    Inwiefern benutzt du damit die Bilinearität des Skalarprodukts
    (s.o.)?
    \Loesung{Mögliche Antworten:
      \begin{gather*}
        \begin{bmatrix}
            0  &-1 & 0\\
            -1 & 4 &-1\\
            0  &-1 & 0
          \end{bmatrix}
        \qquad
        \begin{bmatrix}
            -1 &-1 &-1\\
            -1 & 8 &-1\\
            -1 &-1 &-1
          \end{bmatrix}
      \end{gather*}
    }
  \end{description}
\item \emph{(Bonus: Bildschärfungsfaltung)}
  Betrachte folgenden Kernel, der eine künstliche Schärfung der
  Kanten eines Bildes bewirkt, und verändere die mittlere Zahl nach
  oben und nach unten:
  \begin{center}
    \includegraphics[width=4em]{schaerfen-kern}
  \end{center}
  Was erhälst du bei niedrigeren Werten?
  \Loesung{Das Bild wird bei 4 zur Kantenerkennung, darunter immer
    mehr schwarz.}
  Was passiert bei höheren?
  \Loesung{Das Bild wird immer heller, der Schärfungseffekt
    weniger.}
  \par
  Kannst du die Operation „Bild schärfen” durch die Summe zweier
  Faltungen beschreiben?
  \Loesung{Die Bildschärfungsfaltung addiert eigentlich nur die
    Ergebnisse der Identität, d.h. das ursprüngliche Bild, und
    das Ergebnis einer Kantenerkennung.
    Erhöht man den Wert, wird das ursprüngliche Bild entsprechend
    vielfach aufaddiert und das Resultat somit insgesamt heller.}
\end{enumerate}
\end{document}
